% !TeX root = ../教材题库_全部条目_转写工作稿.tex
% Source: XJ-B1-C06-RV-P267-Y1120
% 第6章 / 章末复习 复习题六 / 教材 p.267、268 / PDF 第274、275页
% OCR未校对：公式与文字均可能有误；未完成数学TeX转写；禁止直接发给学生
% 原题及相邻原解截图（核对用）：../题目原图/XJ-B1-C06-RV-P267-Y1120-01.webp, ../题目原图/XJ-B1-C06-RV-P267-Y1120-02.webp
% 本文件只在转写完成后作为题干使用；OCR 条目尚不是可靠题干。
\begin{bankitem}{复习题六 第14题}{第6章 / 章末复习 复习题六 / 教材 p.267、268 / PDF 第274、275页}{XJ-B1-C06-RV-P267-Y1120}{OCR未校对：公式与文字均可能有误；未完成数学TeX转写；禁止直接发给学生}
\begin{ocrdraft}
[教材p.267 / PDF第274页]
\par
14．为了解一个鱼塘中养殖的鱼的生长情况，从这个鱼塘中多个不同位置捕
\par
捞出100条鱼，称得每条鱼的质量（单位：(g)，并将所得数据分组（每组包含左端
\par
值，不包含右端值），画出频率分布直方图，如图所示．
\par
6
\par
5
\par
4
\par
3
\par
2
\par
1
\par
0
\par
L001．05L10L15L20L25L30质《黾/kg
\par
（第14题）
\par
[教材p.268 / PDF第275页]
\par
（D根据直方图作频率分布表；
\par
（2）估计数据落在[1.巧，L30）中的概率为多少；
\par
（3）将上面捕捞的100条鱼分别做一记号后再放回鱼塘，几天后再从鱼塘的
\par
多处不同位置捕捞出120条鱼，其中带有记号的鱼有6条，请根据这一情况来估
\par
计该鱼塘中鱼的总条数．
\end{ocrdraft}
\end{bankitem}
